% ============================================================
%  12_architecture.tex — Section 12 : Architecture du code
% ============================================================
\section{Architecture du code}

\begin{lstlisting}[style=benzai]
csp_solver/
  main.py                    # Point d'entree CSP, orchestration
                             # Drapeaux : --validate, --method (md=defaut|multi-runs),
                             # --n-runs, --count, --count-hexagon, --md-no-deterministic
  test.py                    # Test sur l'original (toutes tailles a 6)
  utils/
    parser.py                # Lecture Benzai (.graph) -> graphe dual
    table.py                 # Chargement table voisinage (JSON)
    preprocessing.py         # Domaines, filtrage tables, automorphismes
    model.py                 # Modele PyCSP3 + appel ACE
                             # Post-filtre tout-hex (--count-hexagon pour conserver)
    validate.py              # xTB + planarite (test ACP, single-run et multi-runs)
    csp_types.py             # Dataclasses (renomme depuis types.py mai 2026 mi-mois
                             # pour ne pas masquer le module standard `types`)
    validation/              # Strategy Pattern (mai 2026)
      base.py                # ABC ValidationStrategy
      multi_runs.py          # MultiRunsStrategy
      md.py                  # MDStrategy (defaut depuis mai 2026 mi-mois)
      __init__.py            # Registry get_strategy(name, **kwargs)
  reconstruction/
    topology.py              # Phase A : modifications topologiques
    placement.py             # Phase B : placement geometrique BFS
    assembler.py             # Phase C : MolecularGraph + export XYZ
    pipeline.py              # Orchestration reconstruction + validation
                             # Dispatch via get_strategy(method, ...)
                             # Defaut : method="md"
  data/
    table_voisinage.json     # Table regeneree par MD (mai 2026 mi-mois) :
                             # T(5)=63, T(6)=280, T(7)=367 (-41% vs ancienne)
    first.graph              # Test : 3 hexagones lineaires
    second.graph             # Test : 16 hexagones
  experiments/
    batch_main.py            # Lance main.py sur un dossier .graph (1 config)
                             # Defaut : --method md
    batch_all.py             # Lance batch_main.py pour les 8 combinaisons
    batch_test.py            # Test rapide sur les originaux
    count_all.py             # Compte les solutions CSP pures (sans xTB)
                             # 8 configs x N molecules en quelques minutes
                             # [mai 2026 mi-mois]
    aggregate_runs.py        # Agrege les runs multi-N en stats + classification
                             # Reconciliation disque + preserve md_validation existant
    aggregate_md.py          # Agrege les runs MD (additif + reconciliation disque)
                             # [mai 2026 mi-mois : drop des entrees data.json
                             #  obsoletes quand sol_dir disparait]
    view.py                  # Generation des output/hX/view.html
                             # Affichage "Aucune solution CSP" pour mol entierement
                             # gelees (mai 2026 mi-mois)
    update_report.py         # Generation du rapport global (dashboard MD-only)
    bundle.py                # Creation de l'archive portable .zip
    templates/               # HTML/CSS/JS externalises (cf. section 11)
    plane/benzdb/            # Datasets BenzAI DB d'entree (h3, h4, h5, h6, ...)
    output/                  # Resultats des batchs (cf. section 10)
    report/index.html        # Rapport global genere

non_benzenoid_generator/
  main.py                    # Pipeline de generation de la table de voisinage
                             # Depuis mai 2026 mi-mois : optimize utilise MD
                             # (md_then_optimize) au lieu de --opt direct
  core/
    optimizer.py             # Wrapper xtb --opt (legacy, encore utilise par
                             # MultiRunsStrategy)
    optimizer_md.py          # Wrapper xtb --md + --opt
                             # OMP/MKL=1 par defaut pour reduire la variance
  utils/
    planarity.py             # Test ACP (utilise par les 2 pipelines)
    report.py                # Genere rapport_planarite.txt
\end{lstlisting}

\textbf{Flux de données :}

\begin{enumerate}
  \item \texttt{parser.py} lit le \texttt{.graph} $\to$ \texttt{BenzenoidGraph} (dual + patrons)
  \item \texttt{preprocessing.py} calcule domaines, filtre tables, extrait automorphismes
  \item \texttt{model.py} construit le CSP (PyCSP3), compile en XCSP3, appelle ACE
  \item ACE énumère les solutions $\to$ parsing de la sortie XML
  \item \texttt{reconstruction/} reconstruit chaque solution en molécule 3D (XYZ)
  \item \texttt{validate.py} optimise (xTB) et teste la planarité (ACP, seuil 10°)
\end{enumerate}

\textbf{Commandes et drapeaux :}

\begin{center}
\begin{tabular}{lp{9cm}}
  \toprule
  Drapeau & Effet \\
  \midrule
  \texttt{--no-freeze} & Désactive C2 ($b(v) \geq 2$). Les hexagones avec arêtes libres
    séparées deviennent libres (domaine $\{5, 6, 7\}$). \\
  \texttt{--no-table} & Désactive C3 (table de voisinage). Seules C1 et C4 restreignent
    les solutions. \\
  \texttt{--adj-57} & Active C5 (adjacence 5--7). Chaque 5 doit avoir un voisin 7 et
    inversement. \\
  \texttt{--validate} & Reconstruit chaque solution en 3D, optimise avec xTB, et teste
    la planarité. \\
  \texttt{--n-runs N} & Lance \texttt{N} optimisations xTB par solution. Pertinent
    \emph{seulement} si \texttt{-{}-method multi-runs} est aussi passé (depuis mai 2026
    mi-mois). \\
  \texttt{--method M} & \emph{[mai 2026]} Sélectionne la stratégie de validation :
    \texttt{md} (par défaut depuis mai 2026 mi-mois) ou \texttt{multi-runs}.
    Active automatiquement \texttt{aggregate\_md.py} ou \texttt{aggregate\_runs.py}
    en sortie de batch. \\
  \texttt{--md-no-deterministic} & \emph{[mai 2026]} Désactive le forçage
    \texttt{OMP\_NUM\_THREADS=1} + \texttt{MKL\_NUM\_THREADS=1} (pour récupérer
    le parallélisme au prix d'une variance accrue, sub-seuil mais visible). \\
  \texttt{--count} & Affiche seulement le nombre de solutions (à combiner avec
    \texttt{-{}-all} pour vraiment énumérer toutes les solutions). \\
  \texttt{--all} & Force ACE à énumérer toutes les solutions (sans \texttt{-{}-all},
    \texttt{-{}-count} se contente de la première). \\
  \texttt{--count-hexagon} & \emph{[mai 2026 mi-mois]} Conserve la solution
    tout-hexagones (le benzénoïde d'origine) dans la liste. Par défaut, elle est
    \textbf{exclue} car le solveur cible les substitutions non-benzenoïdes. \\
  \bottomrule
\end{tabular}
\end{center}

\begin{lstlisting}[style=benzai]
# === Solveur seul (depuis csp_solver/) ===
python main.py data/fichier.graph                          # Defaut (C1+C2+C3+C4)
python main.py data/fichier.graph --no-freeze              # Sans C2 (b>=2)
python main.py data/fichier.graph --no-table               # Sans C3 (table)
python main.py data/fichier.graph --adj-57                 # Avec C5 (adj 5-7)

# Validation : MD est le defaut depuis mai 2026 mi-mois
python main.py data/fichier.graph --validate               # MD + opt (defaut)
python main.py data/fichier.graph --validate --method multi-runs --n-runs 10
                                                           # Retour multi-runs N=10
python main.py data/fichier.graph --validate --count-hexagon
                                                           # Conserver tout-hex
python main.py data/fichier.graph --count --all            # Compter sans validation

# === Batch sur un dossier (depuis csp_solver/experiments/) ===
python batch_main.py plane/benzdb/h4 --validate            # MD (defaut)
python batch_main.py plane/benzdb/h4 --validate --method multi-runs --n-runs 10
python batch_all.py  plane/benzdb/h4                       # MD sur 8 configs (defaut)
python batch_all.py  plane/benzdb/h4 --method multi-runs --n-runs 10

# Decompte rapide sans xTB (mai 2026 mi-mois)
python count_all.py plane/benzdb/h6                        # 8 configs, ~15 min sur h6

# === Generation des rapports HTML ===
python view.py output/h4 --aggregate                       # viewer h4
python update_report.py                                    # rapport global (dashboard MD)
python bundle.py                                           # archive .zip portable

# === Regeneration de la table de voisinage (depuis non_benzenoid_generator/) ===
python main.py generate                                    # ~1228 XYZ candidats
python main.py optimize                                    # MD + opt [mai 2026 mi-mois]
                                                           # Produit rapport_planarite.txt
# Puis depuis csp_solver/ :
python utils/table.py                                      # Regenere table_voisinage.json
\end{lstlisting}
